% Figure Legends for Manuscript
% Copy and paste into your LaTeX document

\subsection*{Figure 1: GCM Analysis Validates Key Network Switching Findings}

\textbf{A.} Network diagram showing Granger-causal relationships between six brain regions from 310 subjects (fBIRN dataset) at 30\% frequency threshold. Node colors indicate functional networks: red = Central Executive Network (CEN: rPPC, rDLPFC), blue = Salience Network (rFIC, ACC), green = Default Mode Network (DMN: PCC, VMPFC). Edge thickness and opacity scale with frequency across subjects. The dominant bidirectional coupling between rFIC and rDLPFC (58\%, thick black edges) replicates the rFIC-DLPFC connectivity emphasized by Sridharan et al. (2008) while revealing its reciprocal nature in resting-state data. Multiple connections between salience and both CEN and DMN regions confirm rFIC's role in mediating network interactions. \textbf{B.} Quantitative ranking of the 15 strongest Granger-causal connections. Bar colors indicate connection types: purple = Salience$\rightarrow$CEN, orange = Salience$\rightarrow$DMN, red = CEN connections, green = DMN connections. The bidirectional rFIC $\leftrightarrow$ rDLPFC relationship dominates (58.06\%), followed by within-DMN connections (PCC $\leftrightarrow$ VMPFC: $\sim$39\%) and salience-to-DMN pathways (ACC $\rightarrow$ PCC/VMPFC: 36-37\%).

\subsection*{Figure 2: RASL Identifies VMPFC as Primary Structural Driver}

\textbf{A.} Causal structure learning results from RASL applied to the same 310 subjects and brain regions, displayed at 50\% frequency threshold. Network layout and color scheme match Figure 1 for direct comparison. At this moderate threshold, 28 directed edges are present, showing extensive interconnectivity across all three functional networks. The thickest edge (VMPFC $\rightarrow$ rFIC) indicates the dominant structural pathway. \textbf{B.} Top 15 strongest structural causal connections ranked by frequency across 2,496 candidate graph solutions (310 subjects $\times$ $\sim$8 solutions per subject). The VMPFC $\rightarrow$ rFIC connection (89.7\%, green bar) far exceeds all others, emerging as the most consistent structural pathway in the dataset. This finding suggests that the DMN hub (VMPFC) occupies a position upstream of the salience network (rFIC) in the brain's structural causal architecture. Strong connections from VMPFC to multiple target regions (rFIC, PCC, rPPC, rDLPFC) indicate VMPFC's role as a central causal driver across networks.

\subsection*{Figure 3: Complementary Insights from Temporal and Structural Causality}

Side-by-side comparison of network organization revealed by GCM (temporal Granger causality) and RASL (structural causality). Left panel shows GCM network at 50\% threshold, emphasizing the bidirectional rFIC $\leftrightarrow$ rDLPFC temporal coupling (thick curved edges). Right panel shows RASL network at 70\% threshold, where only the unidirectional VMPFC $\rightarrow$ rFIC structural pathway survives this stringent criterion. Bottom panels show comparative bar charts with color-coded connection types. \textit{Similarities:} Both methods identify high cross-network connectivity and central salience network positioning. \textit{Key difference:} GCM emphasizes reciprocal Salience-CEN dynamics (rFIC $\leftrightarrow$ rDLPFC: 58\%), while RASL reveals DMN as structural initiator (VMPFC $\rightarrow$ rFIC: 89.7\%). This complementarity demonstrates that temporal precedence relationships (GCM) and structural causal pathways (RASL) represent distinct organizational principles, both essential for understanding brain network coordination.

\subsection*{Figure 4: Clinical Invariance of Structural Network Architecture}

Comparison of RASL-derived causal networks between healthy controls (Group 0, N=180, left panel) and schizophrenia patients (Group 1, N=130, right panel) at 80\% frequency threshold. Both groups exhibit identical network topology: a single dominant connection from VMPFC to rFIC (Group 0: 89.71\%, Group 1: 89.77\%, difference = 0.06\%). This remarkable similarity across clinical populations indicates that the core structural causal architecture of large-scale brain networks is preserved in schizophrenia. The finding suggests that neuropsychiatric differences may arise from altered network dynamics, connection strength modulation, or task-dependent engagement patterns rather than fundamental reorganization of causal pathways. Node arrangement and colors as in Figures 1-3.

% Supplementary Figures

\subsection*{Figure S1: GCM Networks Across Multiple Thresholds}
GCM-derived networks at four frequency thresholds: 0\% (all edges), 30\%, 40\%, and 50\%. As threshold increases, the network progressively simplifies, with the bidirectional rFIC $\leftrightarrow$ rDLPFC connection consistently surviving as the most robust temporal relationship. At 50\% threshold, only two edges remain (both directions of rFIC-rDLPFC coupling), demonstrating the dominance of this salience-executive link in temporal dynamics. Color scheme as in Figure 1.

\subsection*{Figure S2: RASL Networks Across Multiple Thresholds}
RASL-derived networks at five frequency thresholds: 50\%, 60\%, 70\%, 80\%, and 90\%. Progressive thresholding reveals the hierarchical strength of structural connections. At 70\% and 80\% thresholds, only VMPFC $\rightarrow$ rFIC survives, demonstrating its exceptional robustness as a structural pathway. The sharp transition from 28 edges at 50\% to 1 edge at 70\% indicates a clear separation between the dominant VMPFC $\rightarrow$ rFIC pathway and other connections.

\subsection*{Figure S3: Complete Connectivity Matrices}  
Heatmaps showing complete pairwise connectivity for both GCM (left) and RASL (right) methods. Each cell represents the frequency of a directed edge from source (row) to target (column). Diagonal is masked (no self-connections). Color scale: blue (0\%) to red (100\%). Annotations mark connections exceeding 50\% (GCM) or 70\% (RASL). Matrix view reveals that RASL generally shows higher connectivity levels than GCM (mean: 57\% vs 34\%), but both methods identify similar sets of strong connections despite different primary relationships.

\subsection*{Table S1: Agreement with Sridharan et al. (2008)}
\begin{table}[h]
\centering
\caption{Comparison of key findings between Sridharan et al. (2008) and current GCM implementation}
\begin{tabular}{lcc}
\hline
Finding & Sridharan et al. & Current Study \\
\hline
rFIC central role & \checkmark & \checkmark \\
rFIC $\rightarrow$ DLPFC & \checkmark & \checkmark (58\%) \\
rFIC $\rightarrow$ DMN & \checkmark & \checkmark (30-34\%) \\
ACC $\rightarrow$ DMN & \checkmark & \checkmark (36-37\%) \\
Salience mediates CEN-DMN & \checkmark & \checkmark \\
rFIC bidirectional with DLPFC & Not emphasized & \checkmark (58\%) \\
\hline
\end{tabular}
\end{table}

\subsection*{Table S2: Method Comparison Statistics}
\begin{table}[h]
\centering
\caption{Quantitative comparison of GCM and RASL causal discovery methods}
\begin{tabular}{lcccc}
\hline
Metric & GCM & RASL & Difference & p-value \\
\hline
Mean edge frequency & 34.02\% & 57.04\% & +23.02\% & $<0.001$ \\
Maximum frequency & 58.06\% & 89.74\% & +31.68\% & N/A \\
Edges $>50\%$ & 2 & 28 & +26 & N/A \\
Strongest connection & rFIC$\leftrightarrow$rDLPFC & VMPFC$\rightarrow$rFIC & -- & -- \\
Correlation (30 edges) & -- & -- & $\rho=0.43$ & $<0.05$ \\
\hline
\end{tabular}
\end{table}

% Notes for manuscript integration:
% 1. Replace Figure X, Y, Z, W with your actual figure numbers
% 2. Adjust figure panel letters (A, B, etc.) based on your layout
% 3. Add appropriate journal-specific citations
% 4. Consider moving detailed statistics to supplement
% 5. Adjust language/terminology to match your manuscript style

